\section{Results}

\begin{figure}[t]
    \centering
    \includegraphics[width=\linewidth]{output/step_28_skip_4_1.5_5.5/plots/1_overall_bar.png}
    \caption{Single aggregate score for each of the 10 metrics across all 700 PIE-Bench
    images for KV-Edit (ours). Blue bars indicate higher-is-better metrics (PSNR, SSIM,
    CLIP scores, HPS, Aesthetic Score); red bars indicate lower-is-better metrics
    (Structure Distance, LPIPS, MSE).}
    \label{fig:overall_bar}
\end{figure}

Figure~\ref{fig:overall_bar} reports the aggregate performance of KV-Edit~\cite{zhu2025kvedit}
across all 700 images of PIE-Bench~\cite{piebench2024}.
Background preservation is strong across the board: a PSNR of 37.63\,dB, SSIM of 0.965,
LPIPS of $7.98\times10^{-3}$, and MSE of $3.88\times10^{-4}$ confirm that pixels outside
the edit mask are reconstructed with high fidelity, while a Structure Distance of
$3.80\times10^{-4}$ indicates that global structural coherence is well maintained.
Prompt alignment is competitive at a CLIP Target score of 23.17, though the within-mask
CLIP Target (edit) score of 19.98 is notably lower, reflecting the added difficulty of
aligning a localised region to a text description.
Per-class breakdowns are given in Tables~\ref{tab:bg} and~\ref{tab:clip}.

% ── Comparison with published baselines ──────────────────────────────────────
\subsection{Comparison with Baselines}

Table~\ref{tab:comparison} compares our KV-Edit run against the published baselines
reported in the original paper~\cite{zhu2025kvedit}, using the same metric scaling
(HPS$\times10^2$, LPIPS$\times10^3$, MSE$\times10^4$).
Our implementation achieves \textbf{superior background preservation} relative to the
paper's reported ``Ours'' row: PSNR improves by $+1.76$\,dB (37.63 vs.\ 35.87),
LPIPS improves by $1.94\times10^{-3}$ (7.98 vs.\ 9.92), and MSE improves by
$0.81\times10^{-4}$ (3.88 vs.\ 4.69).
CLIP Target similarity also improves slightly (23.17 vs.\ 22.39).
The HPS\,v2.1 score (24.92) and Aesthetic Score (6.35) are lower than the paper's values
(27.21 and 6.49 respectively), consistent with our observation that the method operates
close to the VAE$^*$ reconstruction regime, faithfully preserving background at the
expense of generating high-perceptual-quality edits within the mask.
Relative to all non-training baselines, KV-Edit achieves the best background preservation
by a substantial margin.

\begin{table}[t]
\centering
\caption{Comparison with published methods on PIE-Bench~\cite{piebench2024}.
VAE$^*$ denotes the inherent reconstruction error through direct VAE reconstruction.
NS indicates no skip step during inversion; RI indicates the reinitialization strategy.
\textbf{Bold} and \underline{underline} denote the best and second-best results.}
\label{tab:comparison}
\setlength{\tabcolsep}{5pt}
\small
\begin{tabular}{lcccccc}
\toprule
\multirow{2}{*}{Method}
  & \multicolumn{2}{c}{Image Quality}
  & \multicolumn{3}{c}{Background Preservation}
  & \multicolumn{1}{c}{Text Align} \\
\cmidrule(lr){2-3}\cmidrule(lr){4-6}\cmidrule(lr){7-7}
  & HPS$_{{\times}10^2}$$\uparrow$
  & AS$\uparrow$
  & PSNR$\uparrow$
  & LPIPS$_{{\times}10^3}$$\downarrow$
  & MSE$_{{\times}10^4}$$\downarrow$
  & CLIP Sim$\uparrow$ \\
\midrule
VAE$^*$     & 24.93 & 6.37 & 37.65 &  7.93 &  3.86 & 19.69 \\
\midrule
P2P~\cite{hertz2022p2p}
            & 25.40 & 6.27 & 17.86 & 208.43 & 219.22 & 22.24 \\
MasaCtrl~\cite{cao2023masactrl}
            & 23.46 & 5.91 & 22.20 & 105.74 &  86.15 & 20.83 \\
RF Inv.~\cite{rfedi2024}
            & \underline{27.99} & \textbf{6.74} & 20.20 & 179.73 & 139.85 & 21.71 \\
RF Edit~\cite{rfedi2024}
            & 27.60 & \underline{6.56} & 24.44 & 113.20 &  56.26 & 22.08 \\
BrushEdit~\cite{brushedit2024}
            & 25.81 & 6.17 & 32.16 &  17.22 &   8.46 & \underline{22.44} \\
FLUX Fill~\cite{bfl2024}
            & 25.76 & 6.31 & 32.53 &  25.59 &   8.55 & 22.40 \\
\midrule
KV-Edit (paper)~\cite{zhu2025kvedit}
            & 27.21 & 6.49 & 35.87 &   9.92 &   4.69 & 22.39 \\
KV-Edit+NS+RI~\cite{zhu2025kvedit}
            & \textbf{28.05} & 6.40 & 33.30 &  14.80 &   7.45 & \textbf{23.62} \\
\midrule
\textbf{KV-Edit (ours)}
            & 24.92 & 6.35 & \textbf{37.63} & \textbf{7.98} & \textbf{3.88} & 23.17 \\
\bottomrule
\end{tabular}
\end{table}

% ── Per-class background preservation ────────────────────────────────────────
\subsection{Per-Class Analysis}

\begin{figure}[t]
    \centering
    \includegraphics[width=\linewidth]{output/step_28_skip_4_1.5_5.5/plots/2_class_heatmap.png}
    \caption{Per-class performance heatmap for KV-Edit (ours) across all edit categories
    and metrics. Green indicates better performance (normalised per metric column);
    red indicates worse performance. Each cell shows the raw metric value.}
    \label{fig:heatmap}
\end{figure}

\begin{table}[t]
\centering
\caption{Per-class background preservation metrics for KV-Edit (ours).
Best value in each column in bold.}
\label{tab:bg}
\setlength{\tabcolsep}{6pt}
\small
\begin{tabular}{lcccc}
\toprule
Edit Class
  & Struct.\ Dist.$\downarrow$
  & PSNR$\uparrow$
  & LPIPS$\downarrow$
  & SSIM$\uparrow$ \\
\midrule
Random              & 0.000414 & 38.387 & 0.00793 & 0.963 \\
Change Object       & 0.000225 & 37.280 & 0.00851 & 0.962 \\
Add Object          & 0.000350 & 38.407 & 0.00701 & 0.972 \\
Delete Object       & \textbf{0.000201} & 35.617 & 0.00989 & 0.956 \\
Change Attr.\ (Content) & 0.000704 & \textbf{38.507} & 0.00766 & \textbf{0.973} \\
Change Attr.\ (Pose)    & 0.000498 & 36.511 & 0.00921 & 0.957 \\
Change Attr.\ (Color)   & 0.000636 & 36.916 & 0.00937 & 0.959 \\
Change Attr.\ (Material)& 0.000234 & 40.818 & \textbf{0.00667} & \textbf{0.973} \\
Change Background   & 0.000409 & 37.080 & 0.00554 & 0.972 \\
\midrule
\textbf{Overall}    & 0.000380 & 37.625 & 0.00798 & 0.965 \\
\bottomrule
\end{tabular}
\end{table}

\begin{table}[t]
\centering
\caption{Per-class CLIP alignment metrics for KV-Edit (ours).
Best value in each column in bold.}
\label{tab:clip}
\setlength{\tabcolsep}{5pt}
\small
\begin{tabular}{lccc}
\toprule
Edit Class
  & CLIP Source$\uparrow$
  & CLIP Target$\uparrow$
  & CLIP Target (edit)$\uparrow$ \\
\midrule
Random              & 25.315 & 22.068 & 20.208 \\
Change Object       & 26.219 & 21.719 & 17.421 \\
Add Object          & 25.239 & 23.941 & 22.263 \\
Delete Object       & 27.110 & 23.194 & 17.039 \\
Change Attr.\ (Content) & 25.241 & 23.911 & 21.369 \\
Change Attr.\ (Pose)    & \textbf{26.710} & \textbf{26.129} & 21.857 \\
Change Attr.\ (Color)   & 26.354 & 23.946 & 19.564 \\
Change Attr.\ (Material)& 25.489 & 23.840 & 22.221 \\
Change Background   & 25.497 & 23.148 & 20.280 \\
\midrule
\textbf{Overall}    & 25.840 & 23.166 & 19.983 \\
\bottomrule
\end{tabular}
\end{table}

\begin{figure}[t]
    \centering
    \begin{subfigure}[b]{0.49\linewidth}
        \includegraphics[width=\linewidth]{output/step_28_skip_4_1.5_5.5/plots/3_distributions.png}
        \caption{Per-image score distributions.}
        \label{fig:boxplot}
    \end{subfigure}
    \hfill
    \begin{subfigure}[b]{0.49\linewidth}
        \includegraphics[width=\linewidth]{output/step_28_skip_4_1.5_5.5/plots/4_radar.png}
        \caption{Normalised radar chart by edit class.}
        \label{fig:radar}
    \end{subfigure}
    \caption{Distribution and per-class radar analysis for KV-Edit (ours).
    \textbf{Left:} box plots showing the spread of per-image scores across all metrics;
    wider boxes indicate higher variance across samples.
    \textbf{Right:} radar chart with per-axis scores normalised to $[0,1]$
    (1 = best across classes); the dashed black line shows the overall mean.}
    \label{fig:dist_radar}
\end{figure}

% ── Discussion ───────────────────────────────────────────────────────────────
\subsection{Discussion}

The KV-Edit baseline establishes a strong floor for background preservation, with the
per-class breakdown revealing several key patterns.
On background preservation (Table~\ref{tab:bg}), Change Attribute (Material) achieves
the highest PSNR (40.82\,dB), indicating that material-level edits—which localise to
surface appearance rather than geometry—minimally disturb surrounding structure.
Delete Object yields the lowest PSNR (35.62\,dB), consistent with the expectation
that object removal requires the model to hallucinate plausible background content,
introducing structural deviation.
Change Attribute (Content) and Material exhibit the highest SSIM ($0.973$), confirming
that attribute edits are the least structurally disruptive edit type.

On text alignment (Table~\ref{tab:clip}), the consistent gap between global CLIP Target
(23.17) and within-mask CLIP Target (edit) (19.98) echoes the pattern observed for
FIA-Edit in the baseline analysis: text-only conditioning provides category-level
guidance but fails to enforce instance-level fidelity within the edited region.
This gap is most pronounced for Delete Object (17.04) and Change Object (17.42),
the two categories that require wholesale replacement of an object with a semantically
different one.
Change Attribute (Pose) achieves the highest CLIP Target score (26.13), indicating that
pose edits are most effectively guided by text descriptions, likely because pose is
a more directly text-encodable attribute than object identity.

Taken together, these results confirm the observation from the baseline FIA-Edit
analysis: background preservation metrics can flatter an editing method's apparent
performance while concealing weak alignment within the edit region.
The within-mask CLIP Target (edit) score of 19.98—three points below the global CLIP
Target—precisely quantifies this limitation and motivates the reference-guided
extension described in Section~3, where visual conditioning replaces text-only
foreground guidance to close this fidelity gap.
